%% 第31回編集委員会（物理第5巻）：組 原W2
%% p5-04.tex  原-4　原子が出す光：エネルギー準位とスペクトル（★2〜4・3段にしない）
%% 中身は計画書 tools/phys5-plan.md の B「原-4」。公式の持ち主は A（A19・A20・A27・A28・A29・A47）。
%%   ほかの章は計画書 B の節の番号で指した（原-2 ③⑥⑪、原-3 ②⑧、原-5、原-7 ③、原-9、原-13、原-1）。
%% 数値はすべて Python で計算した：checks/物理第5巻/原-4/calc_p5_04.py（出力 calc_p5_04_out.txt）。
%% 数学との重なり（N-10・W-6・Y-9）：4-2（円とラジアン）、3-12（対数）、7-8（固有値の予告）は指して使った。
%% 前の巻・数学の本との約束（計画書 A-6・C-6-2）
%%   ・Q13（熱-15:762「レーザーの反転分布は熱のつり合いからはずれた状態（『原子』）」、熱-15 例題の 0.1 eV・300 K で 0.021）→ ⑪・例題11
%%   ・Q17（熱-15 → 原-4 の道しるべ）→ 道しるべの「前に」に熱-15
%%   ・Q33（波-13:105・576「元素は決まった波長の光を出し吸う（原-4）。分光分析」）→ ①〜③
%%   ・Q38（波-13 → 原-4 の道しるべ）→ 道しるべの「前に」に波-13
%%   ・電-15 ⑧（ラーモアの式）→ ⑦（ボーアの模型の限界）
%%   ・波-6（定在波）→ ⑤（円周の定在波）・⑫（レーザーの共振器）
%% 記号（この章）：主量子数 n（カード）、エネルギー準位 E_n、リュードベリ定数 R_∞・R_H（カード。抵抗 R はこの章に出ない）、
%%   ボーア半径 a_0（カード）、クーロンの定数 k（電-1 と同じ。この章で波数を使わない）、原子番号 Z、X線管の電圧 V、
%%   アインシュタイン係数 A・B（カード。⑫だけ。磁場 B とは別）、光の振動数 f、ボルツマン定数 k_B、
%%   1 m^3・1 Hz あたりの光のエネルギー u_f（⑫だけ）、上と下の段の原子の数 N_2・N_1。換算質量は記号を作らず mM/(m+M) と書いた（μ は原-6）。
%% 出典（11月11日に作業環境から確かめた）
%%   ・CODATA 2022（R_∞=10973731.568157 /m、a_0=5.29177210544e-11 m、R_∞hc=13.605693122990 eV、α、m_e、m_p、u）：計画書 D4。
%%   ・原子の線（空気中の波長）：NIST Handbook of Basic Atomic Spectroscopic Data の strong lines の表
%%       （https://physics.nist.gov/PhysRefData/Handbook/Tables/sodiumtable2.htm など。H・Na・Li・K・Ca・Sr・Ba・Cu・Hg・Ne を読んだ）：
%%       Na 588.995・589.592、Li 670.776・670.791、K 766.490、Ca 422.673、Sr 460.733、Ba 553.548、Cu 521.820、Hg 435.833・546.074、
%%       Ne 632.816（He-Ne レーザー）、H 656.285（3本に分かれ 656.271・656.272・656.285）・486.136・434.046・410.174、
%%       ライマン α 121.567（真空）、パッシェン α 1875.101 nm。
%%   ・電離エネルギー：NIST Handbook の表1（H 13.598433・He 24.587387・He⁺ 54.417760・Li 5.391719・Ne 21.56454・Na 5.139076 eV）。
%%   ・ナトリウム 3p の遷移確率 0.616×10⁸ /s：NIST Handbook sodiumtable3（原-3 例題9 と同じ）。
%%   ・特性X線（Kα₁＝KL₃、Kβ＝KM₃）：NIST X-ray Transition Energies（https://physics.nist.gov/cgi-bin/XrayTrans/search.pl）
%%       Fe 6404.006・Cu 8047.823・Mo 17479.372・Ag 22162.917・W 59318.847 eV、Cu Kβ 8905.413 eV。
%%   ・空気の屈折率：Edlén（1966）の標準空気の式で計算（656 nm で 1.000276）。NIST の空気中の波長と 0.003 nm 以内で合う。
%%   ・ノーベル物理学賞（nobelprize.org の summary）：1922 ボーア "for his services in the investigation of the structure of atoms
%%       and of the radiation emanating from them"、1964 タウンズ・バソフ・プロホロフ（メーザー・レーザー）。
%%   ・ボーアのノーベル講演（https://www.nobelprize.org/uploads/2018/06/bohr-lecture.pdf）：1913年の模型、バルマーの式、
%%       リュードベリ、モーズリーの X線の研究。
%%   ・バルマー 1885年・モーズリー 1913〜1914年：二次資料2つ（Britannica の spectral line series、St Andrews の DSB の Balmer の項）
%%       で照らし合わせた（2026年11月）。
%%   ・アインシュタインの A・B 係数（1916〜1917年）：Straumann arXiv:quant-ph/0510180・arXiv:1703.08176 の本文。
%%   ・アンモニアのメーザー（1954年、反転の線 23.870 GHz）：Gordon・Zeiger・Townes, Phys. Rev. 95, 282 (1954)。
%%       二次資料2つ（AIP history "masers-lasers"、Radio-Electronics 1955年6月の記事）で照らし合わせた。
%%   ・花火の色を出すのは主に分子（SrCl・SrOH・BaCl・CuCl）：二次資料2つ（ChemistryViews の記事、Munsell の
%%       "Chemistry of Fireworks Colors"）で照らし合わせた（2026年11月）。分子の帯の波長の範囲は書かなかった。
%%   ・陽子の大きさ：2010年のミューオン水素 0.84184 fm、当時の CODATA 0.8768 fm、水素の精密分光 0.877・0.8482 fm、
%%       CODATA 2022 0.84075 fm〔確〕：arXiv:2501.11195（総説）の本文、計画書 D4。
%%   ・古典の電子が落ちる時間 1.6×10⁻¹¹ s：電-15 ⑧のラーモアの式から計算（calc_p5_04.py）。
%%   ・蛍光灯の水銀の線 435.8・546.1 nm：NIST Handbook〔確〕。
\makeatletter
\tikzset{txe lab/.style={font=\sffamily\footnotesize, text=bkInk}, txe arr/.style={line width=0.9pt, -{Stealth[length=2.2mm]}},
  txe dim/.style={bkGray, line width=0.5pt, {Stealth[length=1.6mm]}-{Stealth[length=1.6mm]}},
  txe lev/.style={bkInk, line width=0.9pt}}
% 目に見える光の帯（380〜780 nm を 0〜120 mm に。#1 は y、#2 は高さ）
\newcommand\txespec[2]{\foreach \w/\r/\g/\b in {380/84/0/88,385/96/0/110,390/105/0/132,395/110/0/155,400/111/0/177,405/108/0/199,410/102/0/222,415/91/0/244,420/74/0/255,425/53/0/255,430/32/0/255,435/11/0/255,440/0/13/255,445/0/38/255,450/0/64/255,455/0/89/255,460/0/115/255,465/0/140/255,470/0/166/255,475/0/191/255,480/0/217/255,485/0/242/255,490/0/255/223,495/0/255/159,500/0/255/96,505/0/255/32,510/9/255/0,515/27/255/0,520/46/255/0,525/64/255/0,530/82/255/0,535/100/255/0,540/118/255/0,545/137/255/0,550/155/255/0,555/173/255/0,560/191/255/0,565/209/255/0,570/228/255/0,575/246/255/0,580/255/245/0,585/255/226/0,590/255/206/0,595/255/186/0,600/255/167/0,605/255/147/0,610/255/128/0,615/255/108/0,620/255/88/0,625/255/69/0,630/255/49/0,635/255/29/0,640/255/10/0,645/255/0/0,650/255/0/0,655/255/0/0,660/255/0/0,665/255/0/0,670/255/0/0,675/255/0/0,680/255/0/0,685/255/0/0,690/255/0/0,695/255/0/0,700/249/0/0,705/238/0/0,710/227/0/0,715/216/0/0,720/205/0/0,725/194/0/0,730/182/0/0,735/171/0/0,740/160/0/0,745/149/0/0,750/138/0/0,755/127/0/0,760/116/0/0,765/104/0/0,770/93/0/0,775/82/0/0}{\definecolor{txetmp}{RGB}{\r,\g,\b}\fill[txetmp] ({(\w-380)*0.22},#1) rectangle ({(\w-375)*0.22+0.05},#1+#2);}}
% 波長 #1（nm）の線を、y=#2 から高さ #3 で（黒い地の上に白い線）
\newcommand\txeline[3]{\fill[white] ({(#1-380)*0.22-0.3},#2) rectangle ({(#1-380)*0.22+0.3},#2+#3);}
% 光子（波の線）：#1 から #2 へ（x の向き）
\newcommand\txephoton[3]{\draw[bkOrange!85!black, line width=0.8pt, -{Stealth[length=2mm]}] plot[domain=#1:#2, samples=60] (\x,{#3+0.9*sin((\x-#1)*90)});}
\makeatother

\phchapter{4}{原子が出す光：\\エネルギー準位とスペクトル}{花火の色は、なぜ材料によって変わる？}{\bklevelrange{2}{4}}{原子が出す光：エネルギー準位とスペクトル}
\bkmichishirube{3-12,4-2,7-8,力-1,電-1,電-15,波-6,波-13,熱-15,原-2,原-3}{原-5,原-7,原-9,原-13}

\begin{goal}
\begin{itemize}[label={\color{bkTheme}●}, leftmargin=1.6\zw, itemsep=0pt, topsep=1pt]
\item \bklevel{2}\enspace 炎色反応と線スペクトルを、原子のエネルギーの段（準位）の差から説明できる。
\item \bklevel{3}\enspace ボーアの模型から水素原子のエネルギー準位を導き、スペクトルの波長を求められる。
\item \bklevel{3}\enspace X線管の電圧から最短波長を、モーズリーの法則から特性X線のエネルギーを求められる。
\item \bklevel{4}\enspace 誘導放出と反転分布、アインシュタインの係数から、レーザーのしくみを説明できる。
\end{itemize}
\end{goal}

\phkoushiki{%
\phk{$mvr=n\dfrac{h}{2\pi}$（量子条件）、$hf=E_n-E_{n'}$（振動数条件）}{高校物理}{ph:原-4:ryoshi}
\phk{$E_n=-\dfrac{13.6}{n^2}$ eV、$r_n=n^2a_0$（水素原子）}{高校物理}{ph:原-4:juni}
\phk{$\dfrac1\lambda=R\Bigl(\dfrac1{n'^2}-\dfrac1{n^2}\Bigr)$（水素のスペクトル）}{高校物理}{ph:原-4:rydberg}
\phk{$\lambda_{\min}=\dfrac{hc}{eV}$（連続X線）、特性X線}{高校物理}{ph:原-4:xsen}
}

\begin{hatena}
夜空に開く花火は、赤・緑・青・黄色と、さまざまな色に光ります。花火の玉には、火薬といっしょに、色を出すための材料が入っています。ナトリウムをふくむ材料なら黄色、ストロンチウムなら赤、バリウムなら緑、銅なら青です。材料を変えるだけで、なぜ、決まった色が出るのでしょう。
\end{hatena}

\noindent 原-2 で、光はエネルギー $hf$ のつぶ（光子）であることを、原-3 で、電子は波でもあることを見ました。この章では、この2つを組み合わせて、原子が決まった色の光だけを出すわけを考えます。答えのかぎは、原子の中の電子が、決まったエネルギーの「段」にしかいられないことです。段の並び方は元素ごとに違うので、出る光の色が元素の「指紋」になります。この章の後半では、同じ考えで、X線と、レーザーのしくみを見ていきます。

\begin{yattemiyou}[分光シートで、いろいろな照明を見比べよう]
\textbf{用意するもの}\par
分光シート（回折格子のシート。理科の教材）、黒い紙のつつ（片方のはしに細いすきまを作り、反対のはしに分光シートをはる）、いろいろな照明（蛍光灯、LED 電球、白熱電球。あれば、ネオンサインやナトリウムランプの街灯）\par\vspace{2pt}
\begin{enumerate}[label={\bkcircled[bkGreen]{\arabic*}}, leftmargin=2.4\zw, itemsep=1pt, topsep=0pt]
\item つつのすきまを照明に向けて、分光シートの側からのぞく。すきまの横に見える色の帯をスケッチする。
\item 白熱電球・LED 電球・蛍光灯で、色の帯がどう違うかを比べる。
\item 蛍光灯で、とくに明るい細い線が何色のところにあるかを記録する。
\end{enumerate}
\textbf{気をつけること}：太陽を直接見ない。理科室の炎色反応の実験は、火と薬品を使うので、自分ではしないで、先生の実験を見る。
\end{yattemiyou}

\noindent やってみようの答え：白熱電球は、赤から紫までとぎれずに続く色の帯（連続スペクトル）に見えます。蛍光灯では、つながった帯の上に、青紫（\qty{436}{nm}）・緑（\qty{546}{nm}）などの、明るく細い線が見えます。これは蛍光灯の中の水銀の原子が出す光です。LED 電球は、青い光と、黄色から赤の幅の広い光の2つの山に見えます（原-9）。\textbf{原子が出す光は、決まった色の細い線になる}――これがこの章の出発点です。

\begin{omoidasou}
\bkrecall{原-2}{光子のエネルギー $E=hf=hc/\lambda$（③）。$hc=\qty{1240}{eV.nm}$ なので、$E$〔eV〕$=1240\div\lambda$〔nm〕。}
\bkrecall{原-3}{電子の波長 $\lambda=h/p=h/(mv)$（②）。}
\bkrecall{波-6}{弦の定在波：長さが半波長の整数倍のときだけ、波が残る（③）。}
\bkrecall{力-1}{円運動の向心力 $\dfrac{mv^2}{r}$。電-1：クーロン力 $\dfrac{ke^2}{r^2}$、$k=\qty{8.99e9}{N.m^2/C^2}$。}
\bkrecall{熱-15}{2つの段の数の比 $e^{-\Delta\varepsilon/(k_{\mathrm B}T)}$（⑦）。0.1 eV・300 K で 0.021。}
\bkrecall{波-13}{回折格子 $d\sin\theta=m\lambda$ で、光を波長ごとに分ける（③）。}
\end{omoidasou}

%% ======================== 物理の道具 ========================
\bksection{doguu}{物理の道具}

\bksubsection{1}{炎色反応と分光}
\label{ph:原-4:sec1}
ガスバーナーの炎に、食塩水をしみこませた針金を入れると、炎があざやかな黄色になります。ほかの物質を入れると、リチウムは赤、カリウムは赤紫、カルシウムはだいだい、銅は青緑と、物質の中の元素によって、決まった色になります。これを\textbf{炎色反応}といいます。\bkidx{えんしょくはんのう}{炎色反応}

色を、もっと詳しく調べるには、光を波長ごとに分けます。プリズムや回折格子（波-13 ③）で光を分ける道具を\textbf{分光器}、分けた光の並びを\textbf{スペクトル}といいます。\bkidx{すぺくとる}{スペクトル}ナトリウムの黄色い炎を分光器で見ると、黄色の所に、細い明るい線がたった2本（\qty{589.0}{nm} と \qty{589.6}{nm}）見えるだけです。

\Needspace{8\baselineskip}\bksubsection{2}{線スペクトルと連続スペクトル}
\label{ph:原-4:sec2}
スペクトルには、大きく3つの形があります（図\ref{ph:原-4:fig-spec}）。
\begin{itemize}[leftmargin=1.2\zw, itemsep=2pt, topsep=2pt, label={\color{bkTheme}●}]
\item \textbf{連続スペクトル}：熱した鉄や白熱電球、太陽のような、熱した固体や濃い気体の光。すべての波長がとぎれずにふくまれ、波長ごとの強さは温度で決まります（熱-13・原-2 ①）。
\item \textbf{線スペクトル}：うすい気体の原子が、放電や炎で光るときの光。決まった波長の細い線だけからなります。\bkidx{せんすぺくとる}{線スペクトル}
\item \textbf{吸収の線}：連続スペクトルの光が、うすい気体を通ると、その気体の原子が出す線と同じ波長の光だけが吸われて、暗い線になります。太陽の光の暗い線は、太陽の外側のうすいガスにある原子が吸ったものです。
\end{itemize}
\textbf{原子は、自分が出す波長の光を、吸うこともできる}――このことも、③で説明します。

\begin{bkfigure}
\begin{tikzpicture}[x=1mm, y=1mm]
  \txespec{30}{7}
  \node[txe lab, anchor=east] at (-2,33.5) {熱した物（連続）};
  \fill[bkInk!90] (0,18) rectangle (88,25);
  \txeline{656.3}{18}{7}\txeline{486.1}{18}{7}\txeline{434.0}{18}{7}\txeline{410.2}{18}{7}
  \node[txe lab, anchor=east] at (-2,21.5) {水素の気体（線）};
  \fill[bkInk!90] (0,6) rectangle (88,13);
  \txeline{589.0}{6}{7}\txeline{589.6}{6}{7}
  \node[txe lab, anchor=east] at (-2,9.5) {ナトリウム（線）};
  \txespec{-6}{7}
  \fill[bkInk] ({(589.0-380)*0.22-0.35},-6) rectangle ({(589.6-380)*0.22+0.35},1);
  \fill[bkInk] ({(656.3-380)*0.22-0.3},-6) rectangle ({(656.3-380)*0.22+0.3},1);
  \fill[bkInk] ({(486.1-380)*0.22-0.3},-6) rectangle ({(486.1-380)*0.22+0.3},1);
  \node[txe lab, anchor=east] at (-2,-2.5) {太陽の光（吸収の線）};
  \foreach \w in {400,450,500,550,600,650,700,750}{\draw[bkInk, line width=0.5pt] ({(\w-380)*0.22},-8) -- ({(\w-380)*0.22},-9.5); \node[txe lab, anchor=north] at ({(\w-380)*0.22},-10) {\w};}
  \node[txe lab, anchor=west] at (90,-11) {nm};
  \foreach \w/\t in {656.3/656, 486.1/486}{\node[font=\sffamily\footnotesize, text=bkInk, anchor=south] at ({(\w-380)*0.22},25.3) {\t};}
  \node[font=\sffamily\footnotesize, text=bkInk, anchor=south east] at ({(410.2-380)*0.22+0.6},25.3) {410};
  \node[font=\sffamily\footnotesize, text=bkInk, anchor=south west] at ({(434.0-380)*0.22-0.6},25.3) {434};
  \node[font=\sffamily\footnotesize, text=bkInk, anchor=south] at ({(589.3-380)*0.22},13.3) {589};
\end{tikzpicture}
\bkfigcaption{いろいろなスペクトル（形をかいた図。横軸は波長）。熱した物はとぎれない色の帯、気体の原子は決まった波長の線（数は空気中の波長 nm）。太陽の光には、ところどころ暗い線がある（ナトリウム・水素の線と同じ波長）。}\label{ph:原-4:fig-spec}
\end{bkfigure}

\Needspace{12\baselineskip}\bksubsection{3}{はてな？の答え：元素の指紋}
\label{ph:原-4:sec3}
\noindent\begin{minipage}[c]{0.55\linewidth}\raggedright
\label{ph:原-4:hatena}\bkhatenaref の答えです。原子の中の電子は、どんなエネルギーでもとれるわけではなく、決まったエネルギーの値（\textbf{エネルギー準位}）にしかいられません。\bkidx{えねるぎーじゅんい}{エネルギー準位}階段の段のように、とびとびです。

炎の熱で上の段（エネルギー $E_2$）に上げられた電子は、やがて下の段（$E_1$）にもどります。そのとき、段の差のエネルギーを、光子1個として出します。光子のエネルギーは $hf$ なので、
\end{minipage}\hfill
\begin{minipage}[c]{0.42\linewidth}\centering
\begin{tikzpicture}[x=1mm, y=1mm]
  \draw[txe lev] (0,0) -- (26,0); \node[txe lab, anchor=west] at (27,0) {$E_1$};
  \draw[txe lev] (0,22) -- (26,22); \node[txe lab, anchor=west] at (27,22) {$E_2$};
  \draw[txe lev, bkGray] (0,30) -- (26,30); \node[txe lab, anchor=west, text=bkGray] at (27,30) {$E_3$};
  \filldraw[fill=bkBlue!25, draw=bkBlue!80!black] (8,22) circle (1.4);
  \draw[txe arr, bkBlue!80!black] (8,20) -- (8,2);
  \txephoton{12}{34}{11}
  \node[txe lab, anchor=south] at (24,13.2) {光子 $hf$};
  \draw[txe dim] (-3,0.4) -- (-3,21.6); \node[txe lab, anchor=east] at (-4,11) {$\Delta E$};
\end{tikzpicture}
\bkfigcaption{上の段から下の段に移ると、差 $\Delta E$ の光子を1個出す。}\label{ph:原-4:fig-dan}
\end{minipage}
\pheq{hf=\frac{hc}{\lambda}=E_2-E_1}{J ＝ J}
\label{ph:原-4:shindou}です。段の差 $\Delta E$ が決まっているので、光の振動数と波長（色）も決まります。反対に、ちょうど $\Delta E$ のエネルギーの光子が来ると、下の段の電子がそれを吸って上の段に上がります（②の吸収の線）。

段の並び方は、原子の中の電子の数と、原子核の電荷で決まるので、元素ごとに違います。だから、出る光の色の組は、元素の「指紋」になります。ナトリウムの黄色い光は、波長 \qty{589}{nm} で、$\Delta E=1240\div589\approx\qty{2.1}{eV}$ の段の差から出ています（例題1）。

\begin{bkcenter}
\begin{bktable}{colspec={Q[l,m,wd=24mm] Q[c,m,wd=26mm] Q[c,m,wd=22mm] Q[c,m,wd=16mm] Q[l,m,wd=24mm]}, colsep=3pt, rowsep=1pt, row{1}={bg=bkTheme!10}}
元素 & 強い線（nm） & 段の差（eV） & 炎の色 & 使われる所 \\
ナトリウム & 589.0・589.6 & 2.11・2.10 & 黄 & 街灯・花火 \\
リチウム & 670.8 & 1.85 & 赤 & 花火 \\
カリウム & 766.5 & 1.62 & 赤紫 & ― \\
カルシウム & 422.7 & 2.93 & だいだい & ― \\
ストロンチウム & 460.7 & 2.69 & 赤 & 花火（赤） \\
バリウム & 553.5 & 2.24 & 黄緑 & 花火（緑） \\
銅 & 510.6・521.8 & 2.43・2.38 & 青緑 & 花火（青） \\
\end{bktable}
\bktabcaption{原子のいちばん強い線（空気中の波長。NIST の表）と段の差（$1240\div$波長）。炎の色が線の色と違う元素があるのは、炎の中では原子でなく分子が光を出すから（せいかくに言うと）。}\label{ph:原-4:tab-sen}
\end{bkcenter}

\begin{seikaku}
花火の赤（ストロンチウム）・緑（バリウム）・青（銅）は、原子の線よりも、炎の中でできる分子（塩化ストロンチウム SrCl・水酸化ストロンチウム SrOH・塩化バリウム BaCl・塩化銅 CuCl など）が出す光が主です。分子にもエネルギーの段がありますが、原子より段が細かく分かれているので、線ではなく幅のある帯になります。表\ref{ph:原-4:tab-sen}のストロンチウムの線（青、460.7 nm）と炎の色（赤）が違うのはこのためです。ナトリウムの黄色やリチウムの赤は、原子そのものの線です。
\end{seikaku}

\begin{phmodel}[とびとびの段をもつ原子]
\phignore{原子が動くことによる波長のずれ（ドップラー効果）、原子どうしの衝突。}
\phassume{原子の中の電子のエネルギーは、決まった段の値だけをとる。段を移るとき、差のエネルギーの光子を1個出すか、吸う。}
\phlimit{分子や固体（段が帯になる：原-8）。とても強い光（光子を2個同時に吸う）。}
\end{phmodel}

\Needspace{8\baselineskip}\bksubsection{4}{水素のスペクトル：バルマーの式}
\label{ph:原-4:sec4}
いちばん簡単な原子は、電子1個の水素です。水素の気体を放電で光らせると、目に見える範囲に4本の線（\qty{656.3}{nm}・\qty{486.1}{nm}・\qty{434.0}{nm}・\qty{410.2}{nm}、空気中の値）が見えます（図\ref{ph:原-4:fig-spec}）。

1885年、スイスの学校の先生だったバルマーは、この4本の波長が、1つの式でぴったり表せることに気づきました。いまの形で書くと、
\pheq{\frac{1}{\lambda}=R\Bigl(\frac{1}{2^2}-\frac{1}{n^2}\Bigr)\qquad(n=3,\ 4,\ 5,\ 6,\ \cdots)}{1/m ＝ 1/m}
です。$R$ は\textbf{リュードベリ定数}とよばれる定数で、約 \qty{1.097e7}{m^{-1}} です。\bkidx{りゅーどべりていすう}{リュードベリ定数}$n=3$ が \qty{656}{nm}、$n=4$ が \qty{486}{nm} …となります（例題3）。

その後、紫外線の中にも、赤外線の中にも、同じような線の並び（\textbf{系列}）が見つかりました。どれも、式の $2^2$ を $1^2$ や $3^2$ に変えるだけで表せます。
\pheq{\frac{1}{\lambda}=R\Bigl(\frac{1}{n'^2}-\frac{1}{n^2}\Bigr)\qquad(n>n')}{1/m ＝ 1/m}
\label{ph:原-4:rydberg}$n'=1$ が\textbf{ライマン系列}（紫外線）、$n'=2$ が\textbf{バルマー系列}（目に見える光）、$n'=3$ が\textbf{パッシェン系列}（赤外線）です。なぜ、こんなに簡単な整数の式になるのか。それを説明したのが、⑤のボーアの模型です。

\begin{phyomi}[$\dfrac1\lambda=R\Bigl(\dfrac{1}{n'^2}-\dfrac{1}{n^2}\Bigr)$]
「波長の逆数は、2つの整数の2乗の逆数の差に比例する」と読みます。$1/\lambda$ は 1 m の中に入る波の数で、光子のエネルギー $hc/\lambda$ に比例します。だから、この式は「光子のエネルギーは、$\dfrac{1}{n'^2}$ と $\dfrac{1}{n^2}$ に比例する2つの量の差」と言いかえられます。これが③の「段の差」です。
\end{phyomi}

\Needspace{8\baselineskip}\bksubsection{5}{ボーアの模型：円周に電子の定在波}
\label{ph:原-4:sec5}
1913年、デンマークのボーアは、水素原子の電子が、原子核（陽子）の周りを円を描いて回っていると考え、2つの条件をおきました。ここでは、ボーアの条件を、原-3 の電子の波で見直します。

\noindent\textbf{条件1：円周に、電子の波がちょうど入る（量子条件）}\par\nobreak
\noindent\begin{minipage}[c]{0.55\linewidth}\raggedright
原-3 で見たように、速さ $v$ の電子は波長 $\lambda=h/(mv)$ の波です。波が円周を1周してもとの所にもどったとき、山と山がそろっていないと、何周もするうちに打ち消し合って消えてしまいます。山がそろうのは、円周 $2\pi r$ が波長のちょうど整数倍のときだけです。弦の定在波（波-6 ③）を、輪にした形です。
\end{minipage}\hfill
\begin{minipage}[c]{0.42\linewidth}\centering
\begin{tikzpicture}[x=1mm, y=1mm]
  \draw[bkGray, line width=0.4pt, dashed] (0,0) circle (11);
  \draw[bkBlue!80!black, line width=1pt, domain=0:360, samples=200, smooth] plot ({(11+2*sin(4*\x))*cos(\x)},{(11+2*sin(4*\x))*sin(\x)});
  \fill[bkRed!70] (0,0) circle (1.2);
  \node[txe lab, anchor=north] at (0,-14.5) {$n=4$：つながる};
  \draw[bkGray, line width=0.4pt, dashed] (32,0) circle (11);
  \draw[bkOrange!85!black, line width=1pt, domain=0:360, samples=200] plot ({32+(11+2*sin(3.5*\x))*cos(\x)},{(11+2*sin(3.5*\x))*sin(\x)});
  \fill[bkRed!70] (32,0) circle (1.2);
  \node[txe lab, anchor=north] at (32,-14.5) {$3.5$ 波長：ずれる};
\end{tikzpicture}
\bkfigcaption{円周の定在波。}\label{ph:原-4:fig-teizai}
\end{minipage}
\[ 2\pi r=n\lambda=n\frac{h}{mv}\qquad(n=1,\ 2,\ 3,\ \cdots) \]
両辺に $mv/(2\pi)$ をかけると、
\pheq{mvr=n\frac{h}{2\pi}=n\hbar}{kg·m/s × m ＝ J·s}
\label{ph:原-4:ryoshi}これを\textbf{量子条件}、整数 $n$ を\textbf{主量子数}といいます。\bkidx{りょうしじょうけん}{量子条件}\bkidx{しゅりょうしすう}{主量子数}

\noindent\textbf{条件2：段を移るとき、差の光子を1個出す（振動数条件）}\par\nobreak
電子が $n$ 番めの段（エネルギー $E_n$）から、低い $n'$ 番めの段に移るとき、
\pheq{hf=E_n-E_{n'}}{J ＝ J}
の光子を1個出します（③と同じ）。これを\textbf{振動数条件}といいます。\bkidx{しんどうすうじょうけん}{振動数条件}

\noindent\textbf{円運動の式と組み合わせる}\par\nobreak
電子（質量 $m$、電荷 $-e$）は、陽子（電荷 $+e$）にクーロン力 $ke^2/r^2$（電-1）で引かれて、半径 $r$ の円を回ります。この力が向心力（力-1）になるので、
\[ \frac{mv^2}{r}=\frac{ke^2}{r^2}\qquad\text{つまり}\qquad mv^2=\frac{ke^2}{r} \tag{ア} \]
量子条件から $v=\dfrac{n\hbar}{mr}$ です。これを（ア）に入れると、
\[ m\Bigl(\frac{n\hbar}{mr}\Bigr)^2=\frac{ke^2}{r}\qquad\text{だから}\qquad r=\frac{\hbar^2}{mke^2}\,n^2 \]
となります。$n=1$ の半径を $a_0$ と書くと、
\pheq{r_n=n^2a_0,\qquad a_0=\frac{\hbar^2}{mke^2}=\qty{5.29e-11}{m}}{m ＝ (J·s)$^2$ ÷ (kg × N·m$^2$/C$^2$ × C$^2$)}
\label{ph:原-4:a0}$a_0$ を\textbf{ボーア半径}といいます（例題5）。\bkidx{ぼーあはんけい}{ボーア半径}

\noindent\textbf{エネルギー}\par\nobreak
電子のエネルギーは、運動エネルギーとクーロン力の位置エネルギー $-ke^2/r$（電-2）の和です。（ア）から $\tfrac12mv^2=\dfrac{ke^2}{2r}$ なので、
\[ E=\frac12mv^2-\frac{ke^2}{r}=\frac{ke^2}{2r}-\frac{ke^2}{r}=-\frac{ke^2}{2r} \]
$r=n^2a_0$ を入れると、
\pheq{E_n=-\frac{ke^2}{2a_0}\times\frac{1}{n^2}=-\frac{13.6\ \mathrm{eV}}{n^2}}{J ＝ N·m$^2$/C$^2$ × C$^2$ ÷ m}
\label{ph:原-4:juni}です。$ke^2/(2a_0)=8.99\times10^9\times(1.60\times10^{-19})^2\div(2\times5.29\times10^{-11})\approx\qty{2.18e-18}{J}$、これを \qty{1.60e-19}{J/eV} でわって \qty{13.6}{eV} です。

\begin{phyomi}[$E_n=-\dfrac{13.6\ \mathrm{eV}}{n^2}$]
「$n$ 番めの段のエネルギーは、$-13.6$ eV を $n^2$ でわった値」と読みます。負の値なのは、電子が陽子から無限に遠くはなれて止まっているときを0にしたからです。$n=1$（いちばん低い段、\textbf{基底状態}）は $-13.6$ eV、$n=2$ は $-3.40$ eV、$n=3$ は $-1.51$ eV と、上の段ほど間がつまり、0 に近づきます。電子を原子から引きはなす（電離する）には、基底状態から 13.6 eV が必要です。\bkidx{きていじょうたい}{基底状態}
\end{phyomi}

\begin{phbaisuru}{主量子数 $n$}
\noindent\begin{minipage}[c]{0.5\linewidth}\raggedright
半径は $n^2$ に比例するので4倍、エネルギーは $1/n^2$ に比例するので $\tfrac14$（0 に近づく）になります。速さは $v=n\hbar/(mr)\propto1/n$ で半分です。$n=1$ の半径 \qty{0.053}{nm} は、$n=2$ で \qty{0.21}{nm}、$n=3$ で \qty{0.48}{nm} です。
\end{minipage}\hfill
\begin{minipage}[c]{0.47\linewidth}\centering
\begin{tikzpicture}[x=1mm, y=1mm]
  \fill[bkRed!70] (0,0) circle (0.8);
  \draw[bkBlue!80!black, line width=0.8pt] (0,0) circle (2);
  \draw[bkOrange!85!black, line width=0.8pt, dashed] (0,0) circle (8);
  \draw[bkPurple!80!black, line width=0.8pt, densely dotted] (0,0) circle (18);
  \node[txe lab, anchor=west] at (20,10) {$n=3$（点線）};
  \node[txe lab, anchor=west] at (20,4) {$n=2$（破線）};
  \node[txe lab, anchor=west] at (20,-2) {$n=1$（実線）};
  \draw[bkGray, line width=0.4pt] (19.5,-2) -- (2.2,-0.5);
\end{tikzpicture}
\end{minipage}
\end{phbaisuru}

\Needspace{12\baselineskip}\begin{phflow}[円運動と量子条件から、水素のスペクトルへ]
\phflowrow{$\dfrac{mv^2}{r}=\dfrac{ke^2}{r^2}$\\$2\pi r=n\dfrac{h}{mv}$}{2つの式から\\$v$ を消す}{$r_n=n^2a_0$\\$E_n=-\dfrac{13.6\ \text{eV}}{n^2}$}
\phflowrow{$hf=E_n-E_{n'}$}{$f=\dfrac{c}{\lambda}$\\$R=\dfrac{13.6\ \text{eV}}{hc}$}{$\dfrac1\lambda=R\Bigl(\dfrac{1}{n'^2}-\dfrac{1}{n^2}\Bigr)$}
\end{phflow}

\noindent\textbf{バルマーの式が出る}\par\nobreak
振動数条件 $hf=hc/\lambda=E_n-E_{n'}$ に $E_n=-13.6\ \text{eV}/n^2$ を入れると、
\[ \frac{hc}{\lambda}=13.6\ \mathrm{eV}\times\Bigl(\frac{1}{n'^2}-\frac{1}{n^2}\Bigr) \]
だから
\[ \frac{1}{\lambda}=\frac{13.6\ \mathrm{eV}}{hc}\Bigl(\frac{1}{n'^2}-\frac{1}{n^2}\Bigr) \]
で、④の式とまったく同じ形になりました。リュードベリ定数は
\[ R=\frac{13.6\ \mathrm{eV}}{hc}=\frac{13.6}{1240\ \mathrm{nm}}=\qty{0.01097}{nm^{-1}}=\qty{1.097e7}{m^{-1}} \]
で、実験で決めた値がぴったり出ます。ボーアは1922年にノーベル物理学賞を受けました。

\begin{seikaku}
ここでは陽子が止まっているとしましたが、本当は、陽子も電子も、2つの重さの中心の周りを回っています。そのため、式の電子の質量 $m$ を、$mM/(m+M)$（$M$ は陽子の質量）にかえるのが正しく、$R$ は $1/(1+m/M)$ 倍（0.99946 倍）になります。CODATA の $R_\infty=\qty{10973731.568}{m^{-1}}$ は「原子核が限りなく重いとき」の値で、水素では $R_{\mathrm H}=\qty{10967758}{m^{-1}}$ です。このわずかな差（約 0.05\,\%）も、スペクトルの精密な測定で確かめられます（例題2）。
\end{seikaku}

\Needspace{14\baselineskip}\bksubsection{6}{エネルギー準位の図と系列}
\label{ph:原-4:sec6}
\begin{bkfigure}
\begin{tikzpicture}[x=1mm, y=1mm]
  % たての位置（−3.4 eV より上は約6倍に広げる）：n=1 0, n=2 25.5, n=3 52.0, n=4 61.2, n=5 65.5, n=6 67.8, ∞ 73.1
  \foreach \y in {0,25.5,52.0,61.2,65.5,67.8}{\draw[txe lev] (0,\y) -- (100,\y);}
  \draw[txe lev, dashed] (0,73.1) -- (100,73.1);
  \foreach \y/\n/\t in {0/1/{$-13.6$}, 25.5/2/{$-3.40$}, 52.0/3/{$-1.51$}, 61.2/4/{$-0.85$}}{
    \node[txe lab, anchor=west] at (101,\y) {$n=\n$};
    \node[txe lab, anchor=east] at (-1,\y) {\t};}
  \node[txe lab, anchor=west] at (101,66.6) {$n=5,6,\cdots$};
  \node[txe lab, anchor=west] at (101,73.6) {$n=\infty$};
  \node[txe lab, anchor=east] at (-1,73.6) {0};
  \draw[bkGray, line width=0.5pt] (-4,11) -- (-2,12.5) -- (-6,13.5) -- (-4,15);
  \node[txe lab, anchor=south, rotate=90] at (-13,38) {エネルギー（eV）};
  % ライマン
  \foreach \x/\y in {8/25.5, 12/52.0, 16/61.2, 20/65.5}{\draw[bkPurple!80!black, line width=0.8pt, -{Stealth[length=1.8mm]}] (\x,\y) -- (\x,0.3);}
  \node[txe lab, anchor=west] at (22,12) {ライマン系列（紫外線）};
  % バルマー
  \foreach \x/\y in {44/52.0, 48/61.2, 52/65.5, 56/67.8}{\draw[bkBlue!80!black, line width=0.8pt, -{Stealth[length=1.8mm]}] (\x,\y) -- (\x,25.8);}
  \node[txe lab, anchor=north] at (50,24.5) {バルマー系列（見える光）};
  % パッシェン
  \foreach \x/\y in {78/61.2, 82/65.5, 86/67.8}{\draw[bkOrange!85!black, line width=0.8pt, -{Stealth[length=1.8mm]}] (\x,\y) -- (\x,52.3);}
  \node[txe lab, anchor=north] at (82,51) {パッシェン系列（赤外線）};
\end{tikzpicture}
\bkfigcaption{水素原子のエネルギー準位と、光を出す移り方。矢印は上の段から下の段への移り方で、矢印が長いほど光子のエネルギーが大きく、波長が短い。ライマン系列は $n=1$ へ、バルマー系列は $n=2$ へ、パッシェン系列は $n=3$ へ移る。見やすいように、$-3.40$ eV より上は、たてを約6倍に広げてかいた（左の折れ線は、目もりの切れ目）。}\label{ph:原-4:fig-juni}
\end{bkfigure}

図\ref{ph:原-4:fig-juni}を\textbf{エネルギー準位の図}といいます。たての長さがエネルギーで、上へ行くほど高い段です。原子から光が出るのは、矢印のように上の段から下の段へ移るときです。$n=1$ に移るライマン系列は、段の差が 10.2 eV 以上で、紫外線です（例題4）。$n=2$ に移るバルマー系列のうち、$n=3$〜6 からの4本が目に見えます。$n=3$ に移るパッシェン系列は、段の差が 1.51 eV より小さく、赤外線です。

$n=\infty$ は、電子が原子からはなれて自由になった状態です。$n=1$ から 0 まで、13.6 eV を与えると電子は原子からはなれます。これを\textbf{電離エネルギー}といいます（例題6）。\bkidx{でんりえねるぎー}{電離エネルギー}NIST の表では、水素の電離エネルギーは \qty{13.598}{eV} です。

\Needspace{8\baselineskip}\noindent\textbf{表計算で、すべての線の波長を出す}\par\nobreak
1行めに上の段 $n$（B〜F列に 2〜6）、A列に下の段 $n'$（2行めから 1〜3）を入れ、B2 に \texttt{=1/(0.010968*(1/\$A2\^{}2-1/B\$1\^{}2))} と入れて右と下にコピーすると、波長〔nm〕の表ができます（$0.010968$ は水素の $R_{\mathrm H}$ を 1/nm で書いた値。$n\leqq n'$ の所は使わない）。
\begin{bkcenter}
\bksheet{6}{17}{%
1 & $n'$ ＼ $n$ & 2 & 3 & 4 & 5 & 6 \\
2 & 1 & 121.6 & 102.6 & 97.3 & 95.0 & 93.8 \\
3 & 2 & ― & 656.5 & 486.3 & 434.2 & 410.3 \\
4 & 3 & ― & ― & 1875.6 & 1282.1 & 1094.1 \\}
\bktabcaption{水素の線の波長〔nm〕（真空の値）。目に見える 380〜780 nm に入るのは、3行めの4つだけ。空気中では、可視光で約 0.03\,\% 短くなる（656.5 → 656.3）。}\label{ph:原-4:tab-sheet}
\end{bkcenter}

\Needspace{12\baselineskip}\bksubsection{7}{ボーアの模型の限界}
\label{ph:原-4:sec7}
ボーアの模型は、水素のスペクトルをみごとに説明しました。けれども、大きな問題をかかえていました。

\begin{itemize}[leftmargin=1.2\zw, itemsep=2pt, topsep=2pt, label={\color{bkTheme}●}]
\item \textbf{回る電子は光を出すはず}：円運動する電子は、いつも中心向きに加速度をもっています。電-15 ⑧のラーモアの式によれば、加速度をもつ電荷は電磁波を出して、エネルギーを失います。計算すると、半径 $a_0$ を回る電子は、約 $1.6\times10^{-11}$ 秒で、うずを巻いて原子核に落ちてしまいます（20万回ほど回るうちに）。ボーアは「決まった段では光を出さない」と、仮定しただけでした。
\item \textbf{電子が2個以上の原子は合わない}：ヘリウム（電子2個）のスペクトルは、ボーアの考え方では正しく出せません。
\item \textbf{線の明るさがわからない}：どの線がどれだけ明るいか（どの移り方が起きやすいか）は、模型からは出ません。
\end{itemize}
\begin{seikaku}
落ちる時間は、ラーモアの式で出ていく光のエネルギーと、半径が小さくなって下がるエネルギーを合わせて求めます。半径 $a_0$ から落ちきるまでの時間は $a_0^3/(4r_{\mathrm e}^2c)\approx\qty{1.6e-11}{s}$ です（$r_{\mathrm e}=\qty{2.82e-15}{m}$ は電子の古典半径とよばれる長さ）。はじめの1周は約 \qty{1.5e-16}{s} なので、そのままわると約10万回ですが、半径が小さくなるほど1周の時間は短くなる（半径の $\tfrac32$ 乗に比例する）ので、回る回数は約2倍の約20万回になります。
\end{seikaku}
これらの答えは、電子を「円を回るつぶ」ではなく、原子核の周りに広がった「波（電子の雲）」として扱う量子力学（原-5・原-7）で得られます。そこでは、電子は加速度をもって回っていないので、光を出しません。そして、おどろいたことに、エネルギー準位は、ボーアの模型と同じ $-13.6\ \text{eV}/n^2$ になります（原-7 ③）。ボーアの模型は、正しいエネルギーを、正しくない絵で出した、すばらしい「近道」だったのです。

\begin{phmodel}[ボーアの模型]
\phignore{原子核の運動（換算質量で直せる：⑤のせいかくに言うと）。電子の速さによる相対論の効果（$v/c\approx1/137$）。}
\phassume{電子は円軌道を回り、円周は電子の波長の整数倍。決まった軌道では光を出さない。}
\phlimit{電子が2個以上の原子。線の明るさ。磁場の中の線の分かれ方（原-7 のスピン）。}
\end{phmodel}

\bksubsection{8}{水素に似たイオン}
\label{ph:原-4:sec8}
ヘリウムの原子から電子を1個取ったイオン He$^{+}$ や、リチウムから2個取った Li$^{2+}$ は、電子を1個だけもっています。原子核の電荷が $+Ze$（$Z$ は原子番号）になるだけなので、⑤のクーロン力を $kZe^2/r^2$ にかえれば、同じ計算ができます。$ke^2$ が $Zke^2$ になるので、
\pheq{r_n=\frac{n^2a_0}{Z},\qquad E_n=-\frac{13.6\ \mathrm{eV}\times Z^2}{n^2}}{m、eV}
\label{ph:原-4:z}です。

\begin{phbaisuru}{原子核の電荷 $Z$}
原子核が電子を強く引くので、半径は $\tfrac12$ に小さくなり、エネルギーは $Z^2=4$ 倍の深さ（$-54.4$ eV）になります。出る光の光子のエネルギーも4倍で、波長は $\tfrac14$ です。NIST の表の He$^{+}$ の電離エネルギーは \qty{54.418}{eV} で、この式とよく合います（例題7）。
\end{phbaisuru}

\Needspace{12\baselineskip}\bksubsection{9}{連続X線と最短波長}
\label{ph:原-4:sec9}
病院の X線写真は、X線管という真空の管で作ります。管の中で、電子を数十 kV の電圧で加速し、タングステンなどの金属の板（ターゲット）にぶつけます。電子が金属の中で急にブレーキをかけられると、光（X線）を出します（電-15 ⑧：加速度をもつ電荷は光を出す）。このX線は、いろいろな波長をふくむ連続スペクトルで、\textbf{連続X線}といいます。\bkidx{れんぞくえっくすせん}{連続X線}

\noindent\begin{minipage}[c]{0.46\linewidth}\raggedright
連続X線には、ある波長より短い光がありません（図\ref{ph:原-4:fig-xsen}）。電子1個が電圧 $V$ で得たエネルギー $eV$ を、全部まとめて1個の光子にしたとき、いちばん短い波長の光子になるからです。$hc/\lambda_{\min}=eV$ から、
\end{minipage}\hfill
\begin{minipage}[c]{0.51\linewidth}\centering
\begin{tikzpicture}
\begin{axis}[bkplot, axis lines=left, width=54mm, height=44mm, xmin=20, xmax=200, ymin=0, ymax=2.25, xtick={41.3,62,100,154}, xticklabels={41,62,100,154}, ytick=\empty, xlabel={波長（pm）}, ylabel={X線の強さ}]
  \addplot[bkcurve, domain=41.33:200, samples=120] {6.75*41.33^3*(1/41.33-1/x)/x^2};
  \addplot[bkcurve2, dashed, domain=62.0:200, samples=120] {6.75*41.33^3*(1/62.0-1/x)/x^2};
  \draw[bkBlue!80!black, line width=1pt] (axis cs:152.8,0.36) -- (axis cs:154.1,2.0) -- (axis cs:155.4,0.35);
  \draw[bkBlue!80!black, line width=1pt] (axis cs:137.9,0.42) -- (axis cs:139.2,0.95) -- (axis cs:140.5,0.41);
  \node[font=\sffamily\footnotesize, anchor=west] at (axis cs:157,1.95) {K$\alpha$};
  \node[font=\sffamily\footnotesize, anchor=east] at (axis cs:136,1.0) {K$\beta$};
\end{axis}
\end{tikzpicture}
\bkfigcaption{銅のターゲットの X線（形をかいた図）。実線は 30 kV、破線は 20 kV。20 kV でも同じ波長に K$\alpha$・K$\beta$ が出る（図では省いた）。}\label{ph:原-4:fig-xsen}
\end{minipage}
\pheq{\lambda_{\min}=\frac{hc}{eV}}{m ＝ J·s × m/s ÷ (C × V)}
\label{ph:原-4:xsen}です。$hc=\qty{1240}{eV.nm}$ を使うと、$\lambda_{\min}$〔pm〕$=1240\div V$〔kV〕です。30 kV で約 \qty{41}{pm}、100 kV で約 \qty{12}{pm} です（例題8）。

\begin{phbaisuru}{X線管の電圧 $V$}
最短波長は $\tfrac12$ になり、もっと短い波長（エネルギーの大きい、物をよく通りぬける）X線が出ます。電圧を上げると、X線の量（強さの合計）も増えます。20 kV の破線と 30 kV の実線を比べると、最短波長が 62 pm から 41 pm に移り、全体が大きくなっています。体の厚い所を写すときや、CT では、100 kV をこえる電圧を使います。
\end{phbaisuru}

\begin{column}{安全：X線を使う機械}
X線は目に見えず、体を通りぬけます。病院の X線写真や空港の手荷物検査の機械は、決まった資格のある人が、X線が外にもれないようにして使います。X線を出す装置を自分で作ったり、動かしたりしてはいけません。体にあたる量（線量）と体への影響は、原-13・原-14 で学びます。
\end{column}

\Needspace{14\baselineskip}\bksubsection{10}{特性X線とモーズリーの法則}
\label{ph:原-4:sec10}
図\ref{ph:原-4:fig-xsen}の連続X線の上には、とがった山がいくつかあります。山の波長は、電圧を変えても動かず、ターゲットの金属の種類で決まります。これを\textbf{特性X線}といいます。\bkidx{とくせいえっくすせん}{特性X線}

\noindent\textbf{内側の段に穴があく}\par\nobreak
重い原子には、電子がたくさんあり、原子核に近い段から順に、電子がうまっています（席の決まりは原-7）。速い電子がぶつかって、いちばん内側（$n=1$）の電子をはじき出すと、そこに穴があきます。すぐに $n=2$ の電子がその穴に落ち、段の差の光子を出します。これが K$\alpha$ 線です。$n=3$ から落ちるのが K$\beta$ 線です。

$n=2$ から落ちる電子から見ると、原子核の電荷 $+Ze$ のうち、$n=1$ に残っているもう1個の電子が $-e$ 分を打ち消すので、およそ $+(Z-1)e$ に引かれていると考えられます。⑧の式で $Z$ を $Z-1$ にすると、
\pheq{E_{\mathrm K\alpha}=13.6\ \mathrm{eV}\times(Z-1)^2\Bigl(\frac{1}{1^2}-\frac{1}{2^2}\Bigr)=10.2\ \mathrm{eV}\times(Z-1)^2}{eV}
\label{ph:原-4:moseley}です。光子のエネルギー（振動数）の平方根が $Z-1$ に比例するこのきまりを、\textbf{モーズリーの法則}といいます。\bkidx{もーずりーのほうそく}{モーズリーの法則}

\noindent\begin{minipage}[c]{0.48\linewidth}\raggedright
イギリスのモーズリーは、1913〜1914年に、いろいろな金属の特性X線をはかり、振動数の平方根が、周期表の順番に、きれいにそろって増えることを見つけました。それまで、周期表の順番（原子番号）は、性質の似た元素を並べた「番号」にすぎませんでした。モーズリーの法則で、原子番号が\textbf{原子核の電荷}そのものであることがわかり、まだ見つかっていない元素の場所（番号の空き）もわかりました。

図\ref{ph:原-4:fig-moseley}で、銅や鉄ではこの式が1\,\%以内で合います。重いタングステンで8\,\%ずれるのは、内側の電子がとても速く（光の速さの半分ほど）、相対論の効果が大きいからです。
\end{minipage}\hfill
\begin{minipage}[c]{0.49\linewidth}\centering
\begin{tikzpicture}
\begin{axis}[bkplot, axis lines=left, width=50mm, height=50mm, xmin=0, xmax=82, ymin=0, ymax=82, xtick={20,40,60,80}, ytick={20,40,60,80}, xlabel={原子番号 $Z$}, ylabel={$\sqrt{E_{\mathrm K\alpha}\div\qty{10.2}{eV}}$}]
  \addplot[bkcurve, domain=1:82] {x-1};
  \foreach \z/\s in {26/25.05, 29/28.08, 42/41.39, 47/46.60, 74/76.24}{\edef\temp{\noexpand\fill[bkOrange!85!black] (axis cs:\z,\s) circle (2.2pt);}\temp}
  \node[font=\sffamily\footnotesize, anchor=north west] at (axis cs:30,26) {Fe・Cu};
  \node[font=\sffamily\footnotesize, anchor=north west] at (axis cs:46,40) {Mo・Ag};
  \node[font=\sffamily\footnotesize, anchor=south east] at (axis cs:73,77) {W};
  \node[font=\sffamily\footnotesize, anchor=north west] at (axis cs:58,52) {$Z-1$（実線）};
\end{axis}
\end{tikzpicture}
\bkfigcaption{モーズリーの法則。実線はこの章の式、点は NIST の測定値（K$\alpha_1$）。}\label{ph:原-4:fig-moseley}
\end{minipage}

\begin{phmijika}[蛍光X線で、物の中の元素を調べる]
物に X線を当てると、物の中の原子の内側の電子がはじき出され、その元素の特性X線（蛍光X線）が出てきます。そのエネルギーをはかれば、どの元素が入っているかが、物をこわさずにわかります。食品にまじった金属の検査、古い絵の絵の具の調査、リサイクルする金属の仕分けなどに使われています。
\end{phmijika}

\Needspace{14\baselineskip}\bksubsection{11}{吸収・自然放出・誘導放出}
\label{ph:原-4:sec11}
ここからは、2つの段（下の段 1、上の段 2、差 $\Delta E=hf$）だけをもつ原子がたくさんあるとして、光とのやりとりを考えます。光とのやりとりには、次の3つがあります（図\ref{ph:原-4:fig-3kotei}）。

\begin{bkfigure}
\begin{tikzpicture}[x=1mm, y=1mm]
  \foreach \X/\t in {12/吸収, 52/自然放出, 92/誘導放出}{
    \draw[txe lev] (\X,0) -- (\X+26,0); \draw[txe lev] (\X,16) -- (\X+26,16);
    \node[txe lab, anchor=east] at (\X-1,0) {1}; \node[txe lab, anchor=east] at (\X-1,16) {2};
    \node[txe lab, anchor=north, font=\sffamily\small\bfseries] at (\X+13,-3) {\t};}
  % 吸収
  \txephoton{0}{17}{8}
  \filldraw[fill=bkBlue!25, draw=bkBlue!80!black] (23,0) circle (1.3);
  \draw[txe arr, bkBlue!80!black] (23,2) -- (23,14);
  \filldraw[fill=bkGray!30, draw=bkGray] (23,16) circle (1.3);
  % 自然放出
  \filldraw[fill=bkBlue!25, draw=bkBlue!80!black] (60,16) circle (1.3);
  \draw[txe arr, bkBlue!80!black] (60,14) -- (60,2);
  \txephoton{64}{80}{8}
  \node[txe lab, anchor=south] at (70,19) {向きはばらばら};
  % 誘導放出
  \txephoton{82}{97}{10}
  \filldraw[fill=bkBlue!25, draw=bkBlue!80!black] (102,16) circle (1.3);
  \draw[txe arr, bkBlue!80!black] (102,14) -- (102,2);
  \txephoton{106}{124}{11}
  \txephoton{106}{124}{5}
  \node[txe lab, anchor=west] at (125,8) {2個};
\end{tikzpicture}
\bkfigcaption{光と原子の3つのやりとり。吸収：光子1個を吸って上へ。自然放出：上の段の原子が、ひとりでに光子を出して下へ。誘導放出：やって来た光子にうながされて、同じ光子をもう1個出して下へ（光子が2個になる）。}\label{ph:原-4:fig-3kotei}
\end{bkfigure}

\begin{enumerate}[label={\bkcircled[bkGreen]{\arabic*}}, leftmargin=2.4\zw, itemsep=1pt, topsep=2pt]
\item \textbf{吸収}：下の段の原子が、エネルギー $hf$ の光子を1個吸って、上の段に上がる。
\item \textbf{自然放出}：上の段の原子が、ひとりでに光子を1個出して、下の段にもどる。どの向きに出るか、いつ出るかは、ばらばら（原-3 ⑧の寿命）。\bkidx{しぜんほうしゅつ}{自然放出}
\item \textbf{誘導放出}：上の段の原子に、エネルギー $hf$ の光子が来ると、その光子にうながされて、\textbf{同じ向き・同じ波長・山と谷のそろった}光子をもう1個出して、下の段にもどる。光子は2個になる。\bkidx{ゆうどうほうしゅつ}{誘導放出}
\end{enumerate}

\noindent\textbf{熱のつり合いでは、上の段は少ない}\par\nobreak
温度 $T$ の周りと熱のつり合いにあるとき、上の段と下の段の原子の数の比は、熱-15 ⑦のボルツマン因子で決まります（2つの段の状態の数が同じとき）。
\pheq{\frac{N_2}{N_1}=e^{-hf/(k_{\mathrm B}T)}}{単位なし}
\label{ph:原-4:boltz}熱-15 の例題では、段の差 0.1 eV・300 K で 0.021 でした。目に見える光の段の差 1.96 eV（赤い光）では、300 K で約 $10^{-33}$ と、ほとんど0です（例題11）。

光が物を通るとき、光子は下の段の原子に吸われて減り、上の段の原子の誘導放出で増えます。吸収は $N_1$ に、誘導放出は $N_2$ に比例する（⑫）ので、熱のつり合い（$N_2<N_1$）では、いつも吸収が勝ち、光は弱まります。光を強めるには、上の段の原子の数を下より多く（$N_2>N_1$）しなければなりません。この状態を\textbf{反転分布}といいます。\bkidx{はんてんぶんぷ}{反転分布}熱-15 で予告したとおり、反転分布は熱のつり合いからはずれた状態で、温度をいくら上げても作れません（$e^{-hf/(k_{\mathrm B}T)}<1$）。電気の放電や強い光で、原子を上の段に「くみ上げ」て作ります。

\Needspace{14\baselineskip}\bksubsection{12}{アインシュタインの係数とレーザー}
\label{ph:原-4:sec12}
1916〜1917年、アインシュタインは、⑪の3つのやりとりの起きやすさを、次のように表しました。光の強さを、1 m$^3$・1 Hz あたりの光のエネルギー $u_f$ で表します。
\begin{itemize}[leftmargin=1.2\zw, itemsep=2pt, topsep=2pt, label={\color{bkTheme}●}]
\item 1秒に吸収する原子の数：$B\,u_f\,N_1$（光が強いほど多い）
\item 1秒に自然放出する原子の数：$A\,N_2$（光の強さによらない）
\item 1秒に誘導放出する原子の数：$B\,u_f\,N_2$（光が強いほど多い）
\end{itemize}
$A$・$B$ を\textbf{アインシュタインの係数}といいます（この節だけの記号。磁場 $B$ とは別）。吸収と誘導放出の係数が同じ $B$ になることは、次のように確かめられます。

熱のつり合いでは、上がる数と下がる数が等しいので、
\[ B\,u_f\,N_1=A\,N_2+B\,u_f\,N_2 \]
です。$N_2/N_1=e^{-hf/(k_{\mathrm B}T)}$ を使って $u_f$ について解くと、
\[ u_f=\frac{A}{B}\times\frac{1}{e^{hf/(k_{\mathrm B}T)}-1} \]
になります。これは、原-2 ⑪で導いたプランクの式 $u_f=\dfrac{8\pi hf^3}{c^3}\times\dfrac{1}{e^{hf/(k_{\mathrm B}T)}-1}$ と、ちょうど同じ形です（吸収と誘導放出の係数を同じとしたから、分母に「$-1$」が出た）。だから
\pheq{\frac{A}{B}=\frac{8\pi hf^3}{c^3}}{J·s/m$^3$}
\label{ph:原-4:ab}です。誘導放出がなければ、プランクの式は出てきません。アインシュタインは、熱した物の光（原-2）から、誘導放出があるはずだと見ぬいたのです。

\begin{phyomi}[自然放出と誘導放出の比]
熱のつり合いの光の中では、自然放出と誘導放出の数の比は $\dfrac{A\,N_2}{B\,u_f\,N_2}=e^{hf/(k_{\mathrm B}T)}-1$ です。目に見える光（$hf\gg k_{\mathrm B}T$）では、この比はとてつもなく大きく、ふだんの光はほとんど自然放出で出ています。電波（$hf\ll k_{\mathrm B}T$）では、比が1より小さく、誘導放出が主になります（例題12）。最初のメーザー（電波のレーザー）が、1954年に、アンモニアの分子の 23.87 GHz の電波で作られたのはこのためです。
\end{phyomi}

\noindent\begin{minipage}[c]{0.43\linewidth}\raggedright
\textbf{レーザーのしくみ}：反転分布にした原子の集まり（ヘリウム・ネオンの気体など）を、2枚の鏡の間に置きます。自然放出で出た光子のうち、鏡に垂直に進むものだけが、鏡の間を何度も行き来し、そのたびに誘導放出で、同じ向き・同じ波長の光子が増えていきます。鏡の間の光は、弦と同じ定在波（波-6）になります。片方の鏡は光の一部を通すので、そこから、向きと波長と山のそろった強い光（\textbf{レーザー光}）が出てきます。\bkidx{れーざー}{レーザー}
\end{minipage}\hfill
\begin{minipage}[c]{0.54\linewidth}\centering
\begin{tikzpicture}[x=1mm, y=1mm]
  \filldraw[fill=bkOrange!10, draw=bkInk, line width=0.6pt] (6,-6) rectangle (54,6);
  \fill[bkInk] (2,-9) rectangle (4,9);
  \fill[bkInk!50] (56,-9) rectangle (57.2,9);
  \draw[bkRed!80!black, line width=0.8pt, domain=4:56, samples=200] plot (\x,{3.5*sin((\x-4)*45)});
  \draw[bkRed!80!black, line width=0.8pt, densely dotted, domain=4:56, samples=200] plot (\x,{-3.5*sin((\x-4)*45)});
  \draw[bkRed!80!black, line width=1.6pt, -{Stealth[length=2.6mm]}] (58,0) -- (70,0);
  \node[txe lab, anchor=south] at (30,9.5) {反転分布の気体（放電でくみ上げる）};
  \node[txe lab, anchor=north] at (3,-10) {鏡};
  \node[txe lab, anchor=north, align=center] at (57,-10) {一部を\\通す鏡};
  \node[txe lab, anchor=north] at (66,-2) {レーザー光};
\end{tikzpicture}
\bkfigcaption{レーザーのしくみ。鏡の間で光が行き来しながら誘導放出で増え、定在波になる。}\label{ph:原-4:fig-laser}
\end{minipage}
\par\vspace{3pt}
レーザー（laser）は「誘導放出による光の増幅」の英語の頭文字です。1964年、メーザーとレーザーのもとになる研究に、ノーベル物理学賞が贈られました。赤い He-Ne レーザーは、ネオンの原子の 632.8 nm の線を使います（例題10）。

\begin{phflow}[熱のつり合いの式から、自然放出と誘導放出の比へ]
\phflowrow{$B\,u_fN_1=AN_2+B\,u_fN_2$\\$\dfrac{N_2}{N_1}=e^{-hf/(k_{\mathrm B}T)}$}{プランクの式\\（原-2 ⑪）と比べる}{$\dfrac{A}{B}=\dfrac{8\pi hf^3}{c^3}$}
\phflowrow{$\dfrac{A}{B}=\dfrac{8\pi hf^3}{c^3}$}{熱のつり合いの\\光の中で}{$\dfrac{\text{自然放出}}{\text{誘導放出}}=e^{hf/(k_{\mathrm B}T)}-1$}
\end{phflow}

%% ---------------------------------------------------------------
\begin{phmada}[陽子の大きさは、どれだけか]
\label{ph:原-4:mada}
水素のスペクトルは、いまでは十数けたもの精度ではかられ、量子力学の計算と比べられています。計算には、陽子の大きさ（電荷の広がり）がほんの少しだけ効いてくるので、逆に、スペクトルから陽子の大きさを求めることができます。

2010年、電子のかわりに、電子の約200倍重い「ミューオン」を陽子の周りに回らせた水素ではかると、陽子の半径は \qty{0.842}{fm} と出ました。ところが、それまでの水素の精密な分光や、電子を陽子にぶつける実験の値は、約 \qty{0.877}{fm} でした。4\,\% の違いは、測定のあいまいさよりずっと大きく、「陽子半径の謎」とよばれました。その後の新しい測定の多くは小さい値に近く、CODATA の2022年の推奨値は \qty{0.84075}{fm} です。けれども、なぜ前の測定が大きい値になったのかは、すべては説明されておらず、測定の見直しと新しい実験が続いています。
\end{phmada}

\begin{phsaki}{原子物理・分光学・レーザー工学へ}
この章では、水素の段を、円を回る電子の模型で求めました。専門の本では、シュレディンガー方程式を解いて、電子の雲の形と、段の細かな分かれ方（電子のスピンや、相対論による「微細構造」）まで計算します（原-5・原-7）。線の明るさ（どの移り方が起きやすいか）も、電子の雲の形から計算できます。レーザーの専門の本では、3つ以上の段を使った反転分布の作り方、鏡の間の光の定在波（共振器のモード）、光の増幅の式を学びます。X線分析では、結晶による X線の回折（原-8）と蛍光X線で、物の中の原子の並びと種類を調べます。
\end{phsaki}

\begin{taisetsu}
\def\tsl#1{\makebox[7.4\zw][l]{\sffamily\bfseries\small\color{\bkTT}#1}}
\tsl{振動数条件}$hf=\dfrac{hc}{\lambda}=E_n-E_{n'}$（段の差の光子を1個出す・吸う）\\[3pt]
\tsl{量子条件}$2\pi r=n\lambda$、$mvr=n\hbar$（円周に電子の定在波）\\[3pt]
\tsl{水素原子}$r_n=n^2a_0$（$a_0=\qty{0.0529}{nm}$）、$E_n=-\dfrac{13.6\ \text{eV}}{n^2}$\\[3pt]
\tsl{スペクトル}$\dfrac{1}{\lambda}=R\Bigl(\dfrac{1}{n'^2}-\dfrac{1}{n^2}\Bigr)$、$R=\qty{1.097e7}{m^{-1}}$\\[3pt]
\tsl{X線}$\lambda_{\min}=\dfrac{hc}{eV}$、$E_{\mathrm K\alpha}\approx10.2\ \text{eV}\times(Z-1)^2$\\[3pt]
\tsl{レーザー}反転分布 $N_2>N_1$。$\dfrac{A}{B}=\dfrac{8\pi hf^3}{c^3}$
\end{taisetsu}

\begin{kotoba}
\small
\bkword{エネルギー準位}{原子の中の電子がとれる、とびとびのエネルギーの値（段）。}
\bkword{線スペクトル}{決まった波長の細い線だけからなるスペクトル。うすい気体の原子が出す。}
\bkword{基底状態}{いちばん低いエネルギーの段にいる状態。}
\bkword{電離エネルギー}{基底状態の電子を、原子から引きはなすのに要るエネルギー。}
\bkword{特性X線}{内側の段にあいた穴に電子が落ちるときに出る、元素で決まる波長のX線。}
\bkword{反転分布}{上の段の原子の数が、下の段より多い状態。レーザーに必要。}
\end{kotoba}

\bkkigocards{\bkkigocard{$n$}{エヌ}{主量子数。段の番号（1、2、3、…）。原-4・5・7。屈折率・数密度の $n$ とは別}{principal quantum number}{原-4}}
{\bkkigocard{$R_\infty$}{アール・インフィニティ}{リュードベリ定数〔1/m〕。約 \qty{1.097e7}{m^{-1}}。水素は $R_{\mathrm H}$}{Rydberg constant}{原-4}}
\bkkigocards{\bkkigocard{$a_0$}{エー・ゼロ}{ボーア半径〔m〕。約 \qty{0.0529}{nm}}{Bohr radius}{原-4}}
{\bkkigocard{$A$・$B$}{エー・ビー}{アインシュタインの係数（自然放出・誘導放出）。原-4 ⑫だけの記号。原-12 の放射能 $A$・原-11 の結合エネルギー $B$ とは別}{Einstein coefficients}{原-4}}

%% ======================== 例題 ========================
\Needspace{10\baselineskip}\bksection{reidai}{例題でたしかめよう}
\label{ph:原-4:ex1}
\begin{phreidai}[花火の色の段の差（はてな？）]
ナトリウムの黄色い光は波長 \qty{589.0}{nm} と \qty{589.6}{nm} の2本、リチウムの赤い光は \qty{670.8}{nm} です。それぞれの段の差 $\Delta E$ を eV で求めましょう（3けたで書きます）。
\phbamen
炎の中で、上の段に上げられた電子が下の段にもどるときに出す光です（③）。
\phmushi
空気中と真空の波長の違い（0.03\,\%）は無視します。原子の動きによる波長のずれも無視します。
\phzu
図\ref{ph:原-4:fig-dan}。$\lambda$〔nm〕、$hc=\qty{1240}{eV.nm}$。
\phriyu
光子1個のエネルギーが段の差に等しい（振動数条件）ので、$\Delta E=hc/\lambda$。
\phtateru
\[ \Delta E=\frac{hc}{\lambda}=\frac{1240}{\lambda\ [\mathrm{nm}]}\ \mathrm{eV} \]
\phtoku
\[ \frac{1240}{589.0}=\qty{2.105}{eV},\qquad\frac{1240}{589.6}=\qty{2.103}{eV},\qquad\frac{1240}{670.8}=\qty{1.849}{eV} \]
\phtashikame
単位：eV·nm÷nm＝eV。\\
大きさ：目に見える光の光子は 1.6〜3.3 eV で、その中に入っています。ナトリウムの2本の差は 0.002 eV と、とても小さい（段がわずかに2つに分かれている：原-7 のスピン）。\\
極端な場合：波長が長い（赤い）ほど、段の差は小さくなります。
\kotae ナトリウム 約\qty{2.11}{eV} と \qty{2.10}{eV}（差は 0.002 eV）、リチウム 約\qty{1.85}{eV}
\end{phreidai}
\begin{jibun}
ストロンチウムの原子の強い線は \qty{460.7}{nm} です。段の差は何 eV ですか（3けたで書く）。
\bkjans{$1240\div460.7\approx\qty{2.69}{eV}$}
\end{jibun}

\label{ph:原-4:ex2}
\begin{phreidai}[水素の赤い線 H$\alpha$ を精密に計算する]
水素の $n=3$ から $n=2$ への光（H$\alpha$）の波長を、陽子の動きの直しを入れて、精密に求めましょう。使う値は、CODATA の $R_\infty=\qty{1.0973732e7}{m^{-1}}$ と、直しの式 $R_{\mathrm H}=R_\infty\times0.999456$ です。空気の屈折率を 1.000276 として、空気中の波長も求め、NIST の表の \qty{656.285}{nm} と比べます（4けた以上で書きます）。
\phbamen
ボーアの模型の式が、どれだけ正確かを確かめる場面です（⑤）。
\phmushi
段の細かな分かれ方（微細構造：NIST の表でも H$\alpha$ は3本に分かれる）は無視し、いちばん強い線と比べます。
\phzu
$n'=2$、$n=3$。真空の波長 $\lambda_0$、空気中の波長 $\lambda=\lambda_0/1.000276$。
\phriyu
水素のスペクトルの式（⑤）。空気の中では光が遅くなり、振動数は同じなので、波長は屈折率の分だけ短くなります（波-13）。
\phtateru
\[ \frac{1}{\lambda_0}=R_{\mathrm H}\Bigl(\frac{1}{2^2}-\frac{1}{3^2}\Bigr)=R_{\mathrm H}\times\frac{5}{36} \]
\phtoku
\[ R_{\mathrm H}=1.0973732\times10^7\times0.999456\approx\qty{1.096776e7}{m^{-1}} \]
\[ \lambda_0=\frac{36}{5\times1.096776\times10^{7}}=\qty{6.5647e-7}{m}=\qty{656.47}{nm} \]
\[ \lambda=\frac{656.47}{1.000276}=\qty{656.29}{nm} \]
\phtashikame
単位：1/(1/m)＝m。\\
大きさ：NIST の \qty{656.285}{nm} と、0.005 nm（100万分の8）の違いで合います。陽子の動きの直しを入れないと \qty{656.11}{nm}（空気中で 655.93 nm）で、0.36 nm もずれます。\\
極端な場合：陽子が限りなく重ければ、直しは0になります。
\kotae 真空で約\qty{656.47}{nm}、空気中で約\qty{656.29}{nm}（NIST の 656.285 nm と合う）
\end{phreidai}
\begin{jibun}
例題2 と同じ $R_{\mathrm H}$ を使って、$n=4$ から $n=2$ への光（H$\beta$）の真空の波長を求めましょう（4けたで書く）。
\bkjans{$1/\lambda_0=R_{\mathrm H}\times3/16$ で、約\qty{486.3}{nm}}
\end{jibun}

\label{ph:原-4:ex3}
\begin{phreidai}[バルマー系列の4本]
バルマー系列（$n'=2$）の、$n=3$・4・5・6 からの光の波長を、$E_n=-13.6\ \text{eV}/n^2$ と $hc=\qty{1240}{eV.nm}$ で求めましょう。
\phbamen
水素の目に見える4本の線を、エネルギー準位から計算する場面です（⑤⑥）。
\phmushi
陽子の動きの直しと、空気の屈折率は無視します（どちらも 0.05\,\% ほど）。
\phzu
図\ref{ph:原-4:fig-juni}のバルマー系列の矢印。
\phriyu
振動数条件 $hc/\lambda=E_n-E_2$（⑤）。
\phtateru
\[ \lambda=\frac{1240}{13.6\times(1/4-1/n^2)}\ \mathrm{nm} \]
\phtoku
\[ n=3:\ \frac{1240}{13.6\times0.1389}=\qty{656}{nm},\qquad n=4:\ \frac{1240}{13.6\times0.1875}=\qty{486}{nm} \]
\[ n=5:\ \frac{1240}{13.6\times0.21}=\qty{434}{nm},\qquad n=6:\ \frac{1240}{13.6\times0.2222}=\qty{410}{nm} \]
\phtashikame
単位：nm。\\
大きさ：NIST の空気中の波長 656.3・486.1・434.0・410.2 nm と合います。\\
極端な場合：$n\to\infty$ では $1240\div3.40=\qty{365}{nm}$（紫外線）で、系列の線はここに集まります。
\kotae \qty{656}{nm}・\qty{486}{nm}・\qty{434}{nm}・\qty{410}{nm}
\end{phreidai}
\begin{jibun}
バルマー系列の5本め（$n=7$ から $n=2$）の波長を求めましょう。目に見えますか。
\bkjans{$\Delta E=13.6\times(1/4-1/49)\approx\qty{3.12}{eV}$、約\qty{397}{nm}。目に見える光のはしで、ほとんど見えない}
\end{jibun}

\label{ph:原-4:ex4}
\begin{phreidai}[ライマン α とパッシェン α]
水素の $n=2\to1$（ライマン α）と、$n=4\to3$（パッシェン α）の光の波長を求め、それぞれ紫外線・可視光・赤外線のどれかを答えましょう。
\phbamen
バルマー系列以外の系列の光です（⑥）。
\phmushi
陽子の動きの直しは無視します。
\phzu
図\ref{ph:原-4:fig-juni}の、ライマン系列とパッシェン系列のいちばん短い矢印。
\phriyu
振動数条件 $hc/\lambda=E_n-E_{n'}$。
\phtateru
\[ \Delta E=13.6\Bigl(\frac{1}{n'^2}-\frac{1}{n^2}\Bigr)\ \mathrm{eV},\qquad\lambda=\frac{1240}{\Delta E}\ \mathrm{nm} \]
\phtoku
\[ \text{ライマン α}:\ \Delta E=13.6\times\frac34=\qty{10.2}{eV},\qquad\lambda=\frac{1240}{10.2}\approx\qty{122}{nm} \]
\[ \text{パッシェン α}:\ \Delta E=13.6\times\Bigl(\frac19-\frac{1}{16}\Bigr)=\qty{0.661}{eV},\qquad\lambda=\frac{1240}{0.661}\approx\qty{1880}{nm} \]
\phtashikame
単位：nm。\\
大きさ：NIST では、ライマン α が \qty{121.6}{nm}、パッシェン α が \qty{1875}{nm}（空気中）で、合います。\\
極端な場合：$n'$ が大きい系列ほど、段の差が小さく、波長が長くなります。
\kotae ライマン α 約\qty{122}{nm}（紫外線）、パッシェン α 約\qty{1.88}{\micro m}（赤外線）
\end{phreidai}
\begin{jibun}
ライマン β（$n=3\to1$）の波長を求めましょう。
\bkjans{$\Delta E=13.6\times8/9\approx\qty{12.1}{eV}$、$1240\div12.1\approx\qty{103}{nm}$（紫外線）}
\end{jibun}

\label{ph:原-4:ex5}
\begin{phreidai}[ボーア半径と電子の速さ]
⑤の式から、水素のボーア半径 $a_0$ と、$n=1$ の電子の速さ $v_1$ を求め、光の速さと比べましょう（3けたで書きます）。
\phbamen
ボーアの模型で、原子の大きさと電子の速さを見積もる場面です（⑤）。
\phmushi
陽子の動きは無視します（陽子は止まっている）。
\phzu
$\hbar=\qty{1.055e-34}{J.s}$、$m=\qty{9.11e-31}{kg}$、\\
$k=\qty{8.99e9}{N.m^2/C^2}$、$e=\qty{1.60e-19}{C}$。
\phriyu
$a_0=\hbar^2/(mke^2)$（⑤）。速さは量子条件 $mv_1a_0=\hbar$ から。
\phtateru
\[ a_0=\frac{\hbar^2}{mke^2},\qquad v_1=\frac{\hbar}{ma_0} \]
\phtoku
\[ a_0=\frac{(1.055\times10^{-34})^2}{9.11\times10^{-31}\times8.99\times10^{9}\times(1.60\times10^{-19})^2}\approx\qty{5.31e-11}{m} \]
丸めた定数の分だけ、少し大きめに出ます。ここからは CODATA の $a_0=\qty{5.29e-11}{m}$ を使います。
\[ v_1=\frac{1.055\times10^{-34}}{9.11\times10^{-31}\times5.29\times10^{-11}}\approx\qty{2.19e6}{m/s} \]
\[ \frac{v_1}{c}=\frac{2.19\times10^6}{3.00\times10^8}\approx\frac{1}{137} \]
\phtashikame
単位：(J·s)$^2$÷(kg·N·m$^2$)＝m。\\
大きさ：CODATA の $a_0=\qty{0.0529}{nm}$ と、丸めた定数の分（0.4\,\%）しか違いません。原子の大きさ（約 0.1 nm）と合います。速さは光の速さの約137分の1で、相対論はほとんど要りません。\\
極端な場合：$n=2$ では、半径4倍、速さ半分です。
\kotae $a_0\approx\qty{0.053}{nm}$（\qty{5.3e-11}{m}）、$v_1\approx\qty{2.19e6}{m/s}$（光の約 $\tfrac{1}{137}$）
\end{phreidai}
\begin{jibun}
$n=2$ の段の半径と電子の速さを求めましょう。
\bkjans{$4a_0\approx\qty{0.212}{nm}$、$v_1/2\approx\qty{1.09e6}{m/s}$}
\end{jibun}

\label{ph:原-4:ex6}
\begin{phreidai}[水素を電離する光]
基底状態の水素原子を、光を当てて電離するには、何 nm より短い波長の光が要りますか。
\phbamen
光子1個で、電子を原子から引きはなす場面です（⑥）。
\phmushi
飛び出した電子の運動エネルギーは0でよいとします（ちょうど電離するぎりぎり）。
\phzu
$n=1$（$-13.6$ eV）から $n=\infty$（0）へ。
\phriyu
光子1個のエネルギーが電離エネルギー以上なら電離できます（光電効果と同じ考え：原-2）。
\phtateru
\[ \frac{hc}{\lambda}\geqq\qty{13.6}{eV}\qquad\text{だから}\qquad\lambda\leqq\frac{1240}{13.6}\ \mathrm{nm} \]
\phtoku
\[ \lambda\leqq\frac{1240}{13.6}\approx\qty{91.2}{nm} \]
\phtashikame
単位：nm。\\
大きさ：紫外線の中でも、とても波長の短い光です。ライマン系列の線は、この 91.2 nm に集まっていきます。\\
極端な場合：これより短い波長の光で電離すると、余ったエネルギーが電子の運動エネルギーになります。
\kotae \qty{91.2}{nm} より短い光（13.6 eV 以上）
\end{phreidai}
\begin{jibun}
$n=3$ の段にいる電子を電離するには、何 eV のエネルギーが要りますか。何 nm より短い光が要りますか。
\bkjans{$13.6\div9\approx\qty{1.51}{eV}$。$1240\div1.51\approx\qty{820}{nm}$ より短い光}
\end{jibun}

\label{ph:原-4:ex7}
\begin{phreidai}[ヘリウムのイオン He$^{+}$]
He$^{+}$（$Z=2$、電子1個）の基底状態のエネルギーと、$n=4\to3$ の光の波長を求めましょう。
\phbamen
水素に似たイオンの段を、$Z^2$ の式で求める場面です（⑧）。
\phmushi
原子核の動きは無視します。
\phzu
$E_n=-13.6\times Z^2/n^2$ eV、$Z=2$。
\phriyu
原子核の電荷が $Z$ 倍になると、エネルギーは $Z^2$ 倍（⑧）。
\phtateru
\[ E_1=-13.6\times2^2,\qquad\Delta E=13.6\times4\times\Bigl(\frac19-\frac{1}{16}\Bigr),\qquad\lambda=\frac{1240}{\Delta E} \]
\phtoku
\[ E_1=\qty{-54.4}{eV},\qquad\Delta E=54.4\times0.04861=\qty{2.645}{eV} \]
\[ \lambda=\frac{1240}{2.645}\approx\qty{469}{nm} \]
\phtashikame
単位：eV と nm。\\
大きさ：NIST の He$^{+}$ の電離エネルギー 54.42 eV と合います。469 nm は青い光で、熱い星のスペクトルに見られます。\\
極端な場合：$Z=1$ にすると、水素のパッシェン α（1875 nm）にもどります（4倍の波長）。
\kotae $\qty{-54.4}{eV}$、約\qty{469}{nm}
\end{phreidai}
\begin{jibun}
Li$^{2+}$（$Z=3$）の基底状態のエネルギーと、$n=2\to1$ の光の波長を求めましょう。
\bkjans{約\qty{-122}{eV}。$\Delta E=122\times\tfrac34\approx\qty{91.8}{eV}$ で、約\qty{13.5}{nm}}
\end{jibun}

\label{ph:原-4:ex8}
\begin{phreidai}[X線管の最短波長]
X線管の電圧が 30 kV のときと 100 kV のときの、連続X線の最短波長を求めましょう。
\phbamen
加速した電子1個のエネルギーが、全部1個の光子になる場面です（⑨）。
\phmushi
加速する前の電子のエネルギーは0とします。
\phzu
図\ref{ph:原-4:fig-xsen}。$V$〔kV〕、$hc=\qty{1240}{eV.nm}$。
\phriyu
$hc/\lambda_{\min}=eV$（⑨）。$eV$ を eV で書くと、$V$〔V〕の数がそのまま eV の数。
\phtateru
\[ \lambda_{\min}=\frac{1240\ \mathrm{eV.nm}}{eV} \]
\phtoku
\[ 30\ \mathrm{kV}:\ \frac{1240}{\num{30000}}=\qty{0.0413}{nm}=\qty{41}{pm} \]
\[ 100\ \mathrm{kV}:\ \frac{1240}{\num{100000}}=\qty{0.0124}{nm}=\qty{12}{pm} \]
\phtashikame
単位：eV·nm÷eV＝nm。\\
大きさ：原子の大きさ（約 100 pm）より短い波長です。電圧を約3倍にすると、最短波長は約 $\tfrac13$ になっています。\\
極端な場合：電圧を下げて数 V にすると、最短波長は目に見える光の波長になります。
\kotae 30 kV で約\qty{41}{pm}、100 kV で約\qty{12}{pm}
\end{phreidai}
\begin{jibun}
CT の X線管で 120 kV を使うときの、最短波長を求めましょう。
\bkjans{$1240\div\num{120000}\approx\qty{0.0103}{nm}$（約\qty{10}{pm}）}
\end{jibun}

\label{ph:原-4:ex9}
\begin{phreidai}[銅・モリブデン・タングステンの K$\alpha$ 線]
モーズリーの法則 $E_{\mathrm K\alpha}=10.2\ \text{eV}\times(Z-1)^2$ を使って、銅（$Z=29$）・モリブデン（$Z=42$）・タングステン（$Z=74$）の K$\alpha$ 線のエネルギーと波長を求め、NIST の測定値（銅 8.048 keV、モリブデン 17.48 keV、タングステン 59.32 keV）と比べましょう。
\phbamen
X線分析の機械が、どの元素かを見分ける手がかりです（⑩）。
\phmushi
内側の電子による打ち消しを、ちょうど電子1個分（$Z-1$）とし、相対論の効果を無視します。
\phzu
図\ref{ph:原-4:fig-moseley}。
\phriyu
$n=2\to1$ の段の差を、電荷 $(Z-1)e$ の水素に似たイオンとして求めます（⑧⑩）。
\phtateru
\[ E=10.2\times(Z-1)^2\ \mathrm{eV},\qquad\lambda=\frac{1240}{E}\ \mathrm{nm} \]
\phtoku
\[ \text{銅}:\ 10.2\times28^2=\qty{8.00}{keV},\qquad\lambda=\frac{1240}{8000}=\qty{0.155}{nm}=\qty{155}{pm} \]
\[ \text{モリブデン}:\ 10.2\times41^2=\qty{17.1}{keV},\qquad\lambda\approx\qty{72}{pm} \]
\[ \text{タングステン}:\ 10.2\times73^2=\qty{54.4}{keV},\qquad\lambda\approx\qty{23}{pm} \]
\phtashikame
単位：eV と pm。\\
大きさ：測定値との差は、銅で 0.6\,\%、モリブデンで 2\,\%、タングステンで 8\,\%。重い元素ほどずれるのは、相対論の効果のためです（⑩）。銅の K$\alpha$ の波長（測定 154 pm）は、結晶の X線回折（原-8）でよく使われます。\\
極端な場合：$Z$ が小さい元素では、K$\alpha$ は X線でなく紫外線になります（水素に似た $Z=1$ ではライマン α）。
\kotae 銅 約\qty{8.0}{keV}（155 pm）、モリブデン 約\qty{17}{keV}（72 pm）、タングステン 約\qty{54}{keV}（23 pm）
\end{phreidai}
\begin{jibun}
鉄（$Z=26$）の K$\alpha$ 線のエネルギーと波長を、モーズリーの法則で求めましょう。
\bkjans{$10.2\times25^2\approx\qty{6.38}{keV}$、$1240\div6380\approx\qty{0.194}{nm}$（測定 6.40 keV）}
\end{jibun}

\label{ph:原-4:ex10}
\begin{phreidai}[He-Ne レーザーの光子の数]
赤い He-Ne レーザー（波長 \qty{632.8}{nm}、出力 \qty{1.0}{mW}）が1秒に出す光子の数を求めましょう。
\phbamen
理科の実験や、昔のバーコード読み取りに使われたレーザーです（⑫）。レーザーの光は、目に入るととても危険です。のぞいたり、人に向けたりしてはいけません。
\phmushi
レーザーの光は、すべて同じ波長の光子とします。
\phzu
$\lambda=\qty{632.8}{nm}$、$P=\qty{1.0e-3}{W}$、光子1個のエネルギー $E$。
\phriyu
1秒に出るエネルギー $P$ を、光子1個のエネルギー $hc/\lambda$ でわると、光子の数です（原-2 ③）。
\phtateru
\[ E=\frac{1240}{632.8}\ \mathrm{eV},\qquad N=\frac{P}{E} \]
\phtoku
\[ E=\qty{1.96}{eV}=1.96\times1.60\times10^{-19}=\qty{3.14e-19}{J} \]
\[ N=\frac{1.0\times10^{-3}}{3.14\times10^{-19}}\approx\num{3.2e15}\ \text{個/s} \]
\phtashikame
単位：W÷J＝1/s。\\
大きさ：1 mW の小さな出力でも、1秒に3000兆個の光子です。1個1個は数えきれないほど多いので、ふだんは光のつぶつぶは感じません。\\
極端な場合：同じ 1 mW でも、波長の長い光ほど光子1個のエネルギーが小さく、数は多くなります。
\kotae 約 \num{3.2e15} 個/s
\end{phreidai}
\begin{jibun}
緑のレーザー（波長 \qty{532}{nm}、出力 \qty{5.0}{mW}）が1秒に出す光子の数を求めましょう。
\bkjans{$E=1240\div532\approx\qty{2.33}{eV}$（\qty{3.73e-19}{J}）。約 \num{1.3e16} 個/s}
\end{jibun}

\label{ph:原-4:ex11}
\begin{phreidai}[熱のつり合いでは、上の段はほとんど空]
段の差が 1.96 eV（He-Ne レーザーの赤い光）の2つの段をもつ原子が、300 K の周りと熱のつり合いにあります。上の段と下の段の数の比 $N_2/N_1$ を求め、熱-15 の例題（0.1 eV で 0.021）と比べましょう。
\phbamen
レーザーに要る反転分布が、温めるだけでは作れないことを確かめる場面です（⑪）。
\phmushi
2つの段の状態の数は同じとし、ほかの段は考えません。
\phzu
$\Delta E=\qty{1.96}{eV}$、$k_{\mathrm B}T=\qty{0.0259}{eV}$（300 K）。
\phriyu
熱のつり合いでは、数の比はボルツマン因子（熱-15 ⑦、⑪）。
\phtateru
\[ \frac{N_2}{N_1}=e^{-\Delta E/(k_{\mathrm B}T)} \]
\phtoku
\[ \frac{\Delta E}{k_{\mathrm B}T}=\frac{1.96}{0.0259}\approx76,\qquad e^{-76}\approx\num{1e-33} \]
\phtashikame
単位：eV÷eV で、単位なし。\\
大きさ：0.1 eV では 0.021（50個に1個）でしたが、1.96 eV では $10^{33}$ 個に1個で、上の段はほとんど空です。反転分布（$N_2/N_1>1$）には、ほど遠いのです。\\
極端な場合：温度を限りなく上げても、$e^{-\Delta E/(k_{\mathrm B}T)}$ は1に近づくだけで、1をこえません。
\kotae 約 $10^{-33}$（0.1 eV の 0.021 よりはるかに小さい）
\end{phreidai}
\begin{jibun}
ナトリウムの段の差 2.10 eV で、炎の温度 2000 K と 3000 K のとき、上の段と下の段の数の比を求めましょう（状態の数の違いは考えない）。
\bkjans{2000 K では $k_{\mathrm B}T\approx\qty{0.172}{eV}$ で約 \num{5e-6}。3000 K では $k_{\mathrm B}T\approx\qty{0.259}{eV}$ で約 \num{3e-4}}
\end{jibun}

\label{ph:原-4:ex12}
\begin{phreidai}[自然放出と誘導放出の比]
300 K の熱のつり合いの光の中で、自然放出と誘導放出の数の比 $e^{hf/(k_{\mathrm B}T)}-1$ を、He-Ne レーザーの光（\qty{632.8}{nm}、1.96 eV）と、アンモニアのメーザーの電波（\qty{23.87}{GHz}）について求めましょう。
\phbamen
光と電波で、どちらの放出が主になるかを比べる場面です（⑫）。
\phmushi
2つの段だけを考えます。
\phzu
$k_{\mathrm B}T=\qty{0.0259}{eV}$。メーザーの光子のエネルギー $hf=6.63\times10^{-34}\times23.87\times10^9$。
\phriyu
アインシュタインの係数から、比は $e^{hf/(k_{\mathrm B}T)}-1$（⑫）。
\phtateru
\[ \frac{\text{自然放出}}{\text{誘導放出}}=e^{hf/(k_{\mathrm B}T)}-1 \]
\phtoku
\[ \text{He-Ne}:\ \frac{hf}{k_{\mathrm B}T}=\frac{1.96}{0.0259}\approx76,\qquad e^{76}-1\approx10^{33} \]
\[ \text{メーザー}:\ hf=\qty{1.58e-23}{J}=\qty{9.89e-5}{eV} \]
\[ \frac{hf}{k_{\mathrm B}T}=\frac{9.89\times10^{-5}}{0.0259}=0.00382,\qquad e^{0.00382}-1\approx0.0038 \]
\phtashikame
単位：どれも単位なし。\\
大きさ：目に見える光では自然放出が圧倒的に多く、電波では誘導放出が約260倍多い（$1/0.0038$）ことがわかります。光のレーザーを作るには、強い反転分布と、鏡で光を何度も往復させることが要ります。\\
極端な場合：$hf\ll k_{\mathrm B}T$ では、比は $hf/(k_{\mathrm B}T)$ に近づきます（$e^x-1\approx x$：5-14）。
\kotae He-Ne で約 $10^{33}$、メーザーで約 0.0038
\end{phreidai}
\begin{jibun}
炭酸ガスレーザー（波長 \qty{10.6}{\micro m}）の光について、300 K での比 $e^{hf/(k_{\mathrm B}T)}-1$ を求めましょう。
\bkjans{$hf=1240\div\num{10600}\approx\qty{0.117}{eV}$、$0.117\div0.0259\approx4.52$、$e^{4.52}-1\approx91$}
\end{jibun}

%% ======================== まちがえやすいところ ========================
\Needspace{11\baselineskip}\bksection{tsumazuki}{まちがえやすいところをチェック}
\begin{tsumazuki}{$n$ が大きい段ほど、エネルギーは小さい？}
\bkng{$E_n=-13.6/n^2$ なので、$n$ が大きいほど数が小さく、エネルギーは低い。}
\bkok{$-13.6/n^2$ は負の数。$n$ が大きいほど0に近づき、エネルギーは\textbf{高い}。}
$-13.6$ eV より $-3.40$ eV のほうが大きい（高い）値です。図\ref{ph:原-4:fig-juni}で、上ほど高い段です。
\end{tsumazuki}

\begin{tsumazuki}{エネルギーの値そのものが、光の色を決める？}
\bkng{$n=3$ の段は $-1.51$ eV なので、1.51 eV の光を出す。}
\bkok{出る光のエネルギーは、2つの段の\textbf{差}。$3\to2$ なら $3.40-1.51=\qty{1.89}{eV}$。}
光は段を移るときに出ます。1つの段の値だけでは決まりません（⑤の振動数条件）。
\end{tsumazuki}

\begin{tsumazuki}{1/λ の式で、λ を出し忘れる？}
\bkng{$R(1/4-1/9)=\qty{1.52e6}{m^{-1}}$ が波長。}
\bkok{これは $1/\lambda$。逆数をとって $\lambda=\qty{6.56e-7}{m}=\qty{656}{nm}$。}
単位が 1/m のままなら、まだ逆数をとっていません。単位で確かめましょう。
\end{tsumazuki}

\begin{tsumazuki}{X線管の電圧を上げると、特性X線の波長も短くなる？}
\bkng{電圧を上げると、すべての X線の波長が短くなる。}
\bkok{短くなるのは連続X線の最短波長だけ。特性X線の波長は、金属の種類で決まり、動かない。}
特性X線は、原子の中の段の差で決まるからです（⑩、図\ref{ph:原-4:fig-xsen}）。20 kV でも、同じ波長に K$\alpha$・K$\beta$ が出ます（銅の内側の電子をはじき出すのに要る約 9 keV より、20 keV のほうが大きいから）。
\end{tsumazuki}

%% ======================== 工学ではこう使う ========================
\Needspace{10\baselineskip}\bksection{phgear}{工学ではこう使う}
\begin{phkougaku}[原子の光で、物をはかる・加工する]
\textbf{発光分光分析}：金属の工場では、製品から小さな試料を取り、放電で光らせて、スペクトルの線の明るさから、鉄の中の炭素やクロムなどの量を数十秒ではかります。水道水の中の金属を調べるときも、試料を高温のプラズマの中で光らせて、元素ごとの線をはかります。

\textbf{蛍光X線分析}：X線を当てて出てくる特性X線のエネルギーで、元素を見分けます（⑩）。電子部品に、使ってはいけない鉛やカドミウムが入っていないかを、こわさずに確かめられます。

\textbf{レーザー加工}：レーザー光は、向きがそろっているので、レンズで小さな点にしぼれます。鉄板を切ったり、溶接したり、スマホの部品に小さな穴をあけたりするのに使われています。光ファイバーの通信や、レーザープリンターも、レーザーの光です。

\textbf{原子時計}：セシウム原子の2つの段の差に当たる電波（振動数 9\,192\,631\,770 Hz）の振動を数えて、1秒の長さを決めています（国際単位系の秒の定義）。

\begin{phketa}[原子の光の目安]
\phketaline{目に見える光の光子}{1.6〜3.3 eV（780〜380 nm）}
\phketaline{原子の段の差}{数 eV（外側の電子）、keV（内側の電子：X線）}
\phketaline{水素}{電離 13.6 eV、ボーア半径 0.053 nm、H$\alpha$ 656 nm}
\phketaline{X線}{最短波長〔pm〕$=1240\div$電圧〔kV〕、特性X線 数〜数十 keV}
\phboxhead{答えが桁で合っているか}
eV と nm は、$1240\div$ で行き来します。目に見える光なら 2 eV くらい、X線なら keV と pm になっているかを確かめましょう。
\end{phketa}
\end{phkougaku}
\phstation{原子物理学,分光分析,X線分析,レーザー工学}{エネルギー準位、線スペクトル、リュードベリ定数、ボーア半径、特性X線、モーズリーの法則、誘導放出、反転分布}

%% ======================== 章末問題 ========================
\bkendproblems
\bkproblemgroup{基本}{3}
\bkproblem カルシウムの原子の強い線は \qty{422.7}{nm} です。この光を出す段の差は何 eV ですか。
\bkpans{$1240\div422.7\approx\qty{2.93}{eV}$}
\bkproblem 水素原子の $n=4$ の段のエネルギーと、軌道の半径を求めましょう。
\bkpans{$-13.6\div16=\qty{-0.850}{eV}$、$16\times\qty{0.0529}{nm}\approx\qty{0.846}{nm}$}
\bkproblem 水素の $n=7$ から $n=3$ への光の波長を求めましょう。紫外線・可視光・赤外線のどれですか。
\bkpans{$\Delta E=13.6\times(1/9-1/49)\approx\qty{1.23}{eV}$、約\qty{1.0}{\micro m}（1005 nm）で赤外線}
\bkproblem X線管の電圧が 50 kV のとき、連続X線の最短波長は何 pm ですか。
\bkpans{$1240\div\num{50000}\approx\qty{0.0248}{nm}$（約\qty{25}{pm}）}
\bkproblem He$^{+}$（$Z=2$）の $n=3$ から $n=2$ への光の波長を求めましょう。水素の H$\alpha$（656 nm）の何倍ですか。
\bkpans{$\Delta E=54.4\times(1/4-1/9)\approx\qty{7.56}{eV}$、$1240\div7.56\approx\qty{164}{nm}$（紫外線）。H$\alpha$ の $\tfrac14$ 倍}
\bkproblemgroup{チャレンジ}{4}
\bkproblem 銀（$Z=47$）の K$\alpha$ 線のエネルギーと波長を、モーズリーの法則で求めましょう。NIST の測定値 22.16 keV との違いは何\,\%ですか。
\bkpans{$10.2\times46^2\approx\qty{21.6}{keV}$、約\qty{57}{pm}。測定値より約 2.6\,\% 小さい}
\bkproblem He-Ne レーザーの光（1.96 eV）で、熱のつり合いの光の中の自然放出と誘導放出の数が等しくなるのは、何 K のときですか。
\bkpans{$e^{hf/(k_{\mathrm B}T)}-1=1$ から $hf/(k_{\mathrm B}T)=\ln2$。$T=hf/(k_{\mathrm B}\ln2)$ で、約\qty{3.3e4}{K}（約3万 K）}
\bkproblem 重水素（原子核は陽子と中性子1個ずつで、陽子の約2倍の質量）の H$\alpha$ 線は、ふつうの水素の H$\alpha$（真空で 656.47 nm）より何 nm 短いですか。$R$ は $1/(1+m/M)$ 倍になること（⑤のせいかくに言うと）を使い、$m/M$ は水素で 1/1836、重水素で 1/3670 とします。
\bkpans{波長の比は約 0.99973。重水素は約\qty{656.29}{nm} で、約\qty{0.18}{nm} 短い}

\begin{furikaeri}
\begin{checklist}
\item 炎色反応と線スペクトルを、原子のエネルギーの段（準位）の差から説明できる。\bkfurref{ph:原-4:sec3}
\item ボーアの模型から水素原子のエネルギー準位を導き、スペクトルの波長を求められる。\bkfurref{ph:原-4:sec5}
\item X線管の電圧から最短波長を、モーズリーの法則から特性X線のエネルギーを求められる。\bkfurref{ph:原-4:sec9}
\item 誘導放出と反転分布、アインシュタインの係数から、レーザーのしくみを説明できる。\bkfurref{ph:原-4:sec11}
\end{checklist}
\end{furikaeri}
\bkfukidashi{mina}{花火の色は、原子や分子の中の電子が段を下りるときに出す光の色だったんだ。段の差が元素ごとに違うから、色が指紋になるんだね！}
